\chapter{Probability, Information, and Entropy}

\lectureoverview{Statistical mechanics begins where microscopic law meets
incomplete knowledge. We first make probability precise, then ask how
deterministic motion itself assigns probabilities, how conservation laws
divide state space, and why information-preserving dynamics leads naturally
to entropy. The first and zeroth laws enter only after this structure is in
place.}

\section{Why Statistical Mechanics?}

Statistical mechanics is sometimes described as old physics. That description
misses the point. It is pre-modern, modern, and likely to remain useful in any
post-modern physics that succeeds the theories we know. Particular models
come and go; the second law of thermodynamics has repeatedly survived as a
guidepost. Whenever a proposed description of nature appears to manufacture
order for free, erase information fundamentally, or run a closed system down
to an unnaturally simple state, the second law warns us that something has
been left out.

The subject may seem less glamorous than a newly discovered particle, but it
is at least as deep and far more general. Gases, liquids, solids, radiation,
black holes, information, and quantum fields all require statistical
reasoning. Indeed, the experimental machinery used to discover particles
could not be understood without it. This is why generations of strong
physicists mastered statistical mechanics even when their principal work lay
elsewhere: the subject is useful, but it is also unusually beautiful.

To see why it is needed, begin with the ideal supplied by mechanics. A closed
system is either the whole universe under consideration or a subsystem
isolated well enough that its surroundings do not influence it. If we know
its initial condition and its law of evolution, then classical mechanics
offers complete prediction. Quantum mechanics modifies what can be predicted,
but it too has a precisely defined state and a law that predicts the maximal
information the theory permits.

There is a severe difference between prediction in principle and useful
understanding. A list of the position and velocity of every molecule in this
room would be fantastically long, would change almost immediately, and would
answer few of the questions we actually care about. Exact microscopic laws
can therefore be perfectly predictive and practically useless at the same
time.

Statistical mechanics is probability theory applied to this physical
situation. It becomes necessary when the initial condition is not known with
perfect accuracy, when the law is known only approximately, when the system
is open, or simply when the number of microscopic variables is too large to
track. None of this implies weak prediction. Large numbers often make
statistical predictions extraordinarily sharp.

Consider an isolated box of gas. We cannot predict the trajectory of every
molecule, yet a few macroscopic data allow us to predict pressure, energy,
and many other bulk quantities with high precision. What remains uncertain
are the fluctuations. At an unpredictable moment, a small region may contain
more molecules than average while another contains fewer. Such an event lies
in the tail of a probability distribution: unlikely, but not forbidden. We
can calculate how often such fluctuations should occur without predicting
the instant of a particular occurrence.

Coin tossing displays the same distinction. In a billion fair tosses, the
numbers of heads and tails will be close to equal. A long run of heads will
eventually occur if we toss often enough, but probability theory predicts its
frequency rather than its appointment in advance. Statistical mechanics is
designed for systems whose elements are too small, too numerous, or both.

The course will develop the laws as well as their applications to matter.
The subject rewards repeated study: one learns it, forgets parts of it, and
learns it again with a deeper sense of which ideas are fundamental. We begin
with the idea that controls everything that follows.

\section{Probability as a Primitive Language}

Why probability works is not something the calculus of probability itself
explains. A probabilistic prediction can fail on a particular trial; that is
what it means for an outcome to be possible but unlikely. The theory tells us
how the exceptions are distributed. It does not abolish them. We therefore
take probability as a primitive concept and state the rules it obeys.

\subsection{States and normalization}

Let \(i\) label a possible outcome or a possible state of a system. For one
coin, \(i\) may mean heads or tails. For one die, it may run from \(1\) to
\(6\). For two dice, a pair of indices records both faces. We begin with a
finite discrete state space,
\begin{equation}
 i=1,\ldots,N,
\end{equation}
and later pass to infinitely many and continuously many states.

Ignorance does not mean that the system has no state. It means that we do not
know which state it has, so we assign a probability \(p_i\) to each
possibility. The basic rules are
\begin{align}
 p_i &\geq 0, \label{eq:sm1-prob-positive}\\
 \sum_{i=1}^{N}p_i &=1. \label{eq:sm1-prob-normalization}
\end{align}
The alternatives are mutually exclusive and exhaustive: exactly one of them
occurs. Negative probabilities have no role here, and the sum is one because
some allowed outcome must occur.

\subsection{Frequency and the law of large numbers}

Probability is connected to experiment by repeated trials. Suppose the same
experiment is performed \(N_{\mathrm{trials}}\) times, and let \(N(i)\) be
the number of trials that produce outcome \(i\). Then the frequency
interpretation is
\begin{equation}
 p_i=
 \lim_{N_{\mathrm{trials}}\to\infty}
 \frac{N(i)}{N_{\mathrm{trials}}}.
 \label{eq:sm1-frequency}
\end{equation}
This is a physical hypothesis expressed through an ideal limit. No laboratory
performs infinitely many experiments. The claim is that sufficiently many
trials make the ratio settle close enough to \(p_i\) for the intended
prediction. Finite samples fluctuate around the limiting value.

\subsection{Averages}

Now assign a number \(f(i)\) to each state. It may be energy, momentum, or any
other quantity we wish to measure. Its expectation value is
\begin{equation}
 \langle f\rangle=\sum_i f(i)p_i.
 \label{eq:sm1-expectation}
\end{equation}
Bracket notation is convenient even outside quantum mechanics. In frequency
language the same average is
\begin{equation}
 \langle f\rangle
 =
 \lim_{N_{\mathrm{trials}}\to\infty}
 \frac{1}{N_{\mathrm{trials}}}
 \sum_i f(i)N(i).
 \label{eq:sm1-frequency-average}
\end{equation}

For a coin, define
\begin{equation}
 f(\mathrm{H})=+1,
 \qquad
 f(\mathrm{T})=-1.
\end{equation}
If heads and tails each have probability \(1/2\), then
\begin{equation}
 \langle f\rangle
 =(+1)\frac12+(-1)\frac12=0.
\end{equation}
The average need not be a possible outcome. The coin never lands on zero;
zero summarizes the distribution of its two possible values.

This brief interlude fixes the vocabulary used throughout the subject:
states, probabilities, frequencies, and expectation values.

\section{Three Ways to Assign Probability}

\subsection{Symmetry}

Why do we usually assign equal probabilities to the two faces of a coin?
The first answer is symmetry. Turning a nearly ideal coin over exchanges
heads and tails without changing the experiment. There is then no physical
reason to favor either outcome. A real mark or imperfection introduces a tiny
bias, but the symmetric model captures the intended experiment:
\begin{equation}
 p(\mathrm{H})=p(\mathrm{T})=\frac12.
\end{equation}

The same argument gives probability \(1/6\) to each face of a symmetric die.
More generally, configurations related by a symmetry are assigned equal
probabilities unless some additional information breaks that symmetry.

\subsection{Experiment}

A misshapen die has no symmetry that exchanges all six faces. In that case,
symmetry alone gives no answer. One may possess a deeper mechanical model
that predicts the bias, but absent such a model the honest procedure is
experimental: throw the die many times, record the frequencies, wait for
them to stabilize, and use the resulting table of probabilities in later
predictions.

Symmetry and experiment are not competing definitions. They are two sources
of information about a probability distribution. Statistical mechanics will
rely heavily on symmetry, but dynamics supplies a third route.

\subsection{Probability from deterministic motion}

Imagine an irregular object with six colored configurations. It changes state
at fixed discrete time intervals according to a definite update law:
\begin{equation}
 R\longrightarrow B\longrightarrow Y\longrightarrow G
 \longrightarrow O\longrightarrow P\longrightarrow R.
 \label{eq:sm1-six-cycle}
\end{equation}
The object itself need not possess any geometric symmetry. The law merely
states which configuration follows a given configuration.

\begin{figure}[htbp]
\centering
\includegraphics[width=0.92\textwidth]{lecture_01_figure_02.png}
\caption{The expectation-value formulas remain on the left board while the
six deterministic color updates are listed on the right.}
\lectureframe{00:26:28}
\end{figure}

The board list is more transparent when drawn as a directed cycle:
\begin{figure}[htbp]
\centering
\begin{tikzpicture}[
  >=Latex,
  state/.style={draw,circle,minimum size=8mm,inner sep=1pt}
]
\node[state] (R) at (90:2.2) {$R$};
\node[state] (B) at (30:2.2) {$B$};
\node[state] (Y) at (-30:2.2) {$Y$};
\node[state] (G) at (-90:2.2) {$G$};
\node[state] (O) at (-150:2.2) {$O$};
\node[state] (P) at (150:2.2) {$P$};
\draw[->] (R) -- (B);
\draw[->] (B) -- (Y);
\draw[->] (Y) -- (G);
\draw[->] (G) -- (O);
\draw[->] (O) -- (P);
\draw[->] (P) -- (R);
\end{tikzpicture}
\caption{A single six-state orbit. Every state is visited once during each
period.}
\label{fig:sm1-six-cycle}
\end{figure}

Suppose we do not know the initial state and take a snapshot at an unknown
time. Since each step lasts equally long, the system spends one sixth of a
period in each color:
\begin{equation}
 p(R)=p(B)=p(Y)=p(G)=p(O)=p(P)=\frac16.
 \label{eq:sm1-uniform-cycle}
\end{equation}
The equality does not come from the shape of the object. It comes from equal
time spent in the states.

We can reorder the arrows and obtain a different microscopic law while still
forming one six-state cycle. The detailed sequence changes, but the snapshot
probability remains \(1/6\) because every state still occurs once per period.
Thus this prediction requires neither the starting point nor the detailed
ordering. It requires only the orbit structure and equal time per update.

\begin{classroomqa}
\classquestion{If the object is not symmetric, in what sense is it
legitimate to give every color probability \(1/6\)?}
\lecturerresponse{The assignment is dynamical rather than geometric. A
snapshot taken at an unknown phase of a six-step orbit samples six equal time
slots. The relevant equivalence is among positions on the orbit, not among
faces of the object. If the dwell times were unequal, the probabilities
would be unequal as well.
\lecturetimestamp{00:27:11--00:29:49}}
\end{classroomqa}

\section{Disconnected Cycles and Conserved Quantities}

A deterministic law need not place all states on one orbit. Consider two
disconnected cycles,
\begin{equation}
 R\longrightarrow B\longrightarrow G\longrightarrow R,
 \qquad
 P\longrightarrow Y\longrightarrow O\longrightarrow P.
 \label{eq:sm1-two-cycles}
\end{equation}
If we know that the system lies on the first cycle, then
\begin{equation}
 \begin{aligned}
 p(R)=p(B)=p(G)&=\frac13,\\
 p(P)=p(Y)=p(O)&=0.
 \end{aligned}
\end{equation}
Knowledge that it lies on the second cycle reverses these assignments.

The more general case requires prior weights. Label the cycles \(+1\) and
\(-1\), and let their probabilities be \(p_+\) and \(p_-\), with
\begin{equation}
 p_++p_-=1.
\end{equation}
Then
\begin{align}
 p(R)=p(B)=p(G)&=\frac{p_+}{3},\\
 p(P)=p(Y)=p(O)&=\frac{p_-}{3}.
 \label{eq:sm1-sector-weights}
\end{align}
The factors \(1/3\) follow from motion within a cycle. The numbers \(p_+\)
and \(p_-\) do not. They must come from additional knowledge, perhaps another
experiment or another system that selects the initial sector.

The labels \(+1\) and \(-1\) define a conserved quantity: once the system
starts on one cycle, the update law never carries it to the other. One may
call the quantity energy, or give it a placeholder name such as
\emph{zilch}. Its name is secondary. Conservation means that state space is
partitioned into regions the dynamics cannot cross.

The cycles need not contain the same number of states. A four-state orbit and
a two-state orbit still define two conserved sectors. Equal-time sampling
then gives \(1/4\) within the first and \(1/2\) within the second, multiplied
in each case by the prior probability of that sector. Statistical mechanics
must answer two logically separate questions:
\begin{enumerate}
\item Given a conserved sector, how much time does the dynamics spend in its
states?
\item What probability should be assigned to the conserved sector itself?
\end{enumerate}

\begin{classroomqa}
\audiencequestion{Are we assuming that all states inside one conserved
sector receive the same probability?}
\lecturerresponse{For the discrete toy laws shown here, yes, because each
state lasts for the same time interval and the motion is deterministic. The
ignorance is in the initial phase of the orbit. A detector with millisecond
timing cannot resolve microsecond updates, although a sufficiently fast
snapshot can still identify the state that was caught. In a more general
system, unequal dwell times must be included explicitly.
\lecturetimestamp{00:36:16--00:37:17}}
\end{classroomqa}

\begin{classroomqa}
\audiencequestion{If two snapshots are taken and the first identifies one
cycle, should the second identify the same cycle?}
\lecturerresponse{Yes. That is precisely what conservation means. Once a
measurement determines the value of the conserved quantity, future
probabilities can be conditioned on that value, and states in the other
sector receive probability zero.
\lecturetimestamp{00:37:18--00:37:45}}
\end{classroomqa}

\subsection{Energy shared between two systems}

Now replace the abstract label by energy. One isolated box is a closed system
and has a conserved energy. Two isolated boxes have two separately conserved
energies because they do not interact. Connect them by a narrow passage that
allows energy to flow, and the separate energies can change while the total
remains fixed.

\begin{figure}[htbp]
\centering
\includegraphics[width=0.92\textwidth]{lecture_01_figure_03.png}
\caption{Two subsystems connected by a narrow channel through which energy
can pass.}
\lectureframe{00:39:19}
\end{figure}

\begin{figure}[htbp]
\centering
\begin{tikzpicture}[>=Latex]
\draw[thick,rounded corners=4pt] (-4,1.8) rectangle (-0.7,-1.5);
\draw[thick,rounded corners=4pt] (0.7,1.5) rectangle (4,-1.8);
\draw[thick] (-0.7,0.45) -- (0.7,0.45);
\draw[thick] (-0.7,-0.35) -- (0.7,-0.35);
\draw[<->,very thick] (-0.45,0.05) -- (0.45,0.05);
\node at (-2.35,0.15) {subsystem \(1\)};
\node at (2.35,-0.15) {subsystem \(2\)};
\node[above] at (0,0.5) {energy exchange};
\end{tikzpicture}
\caption{Once the connection is opened, only \(E_1+E_2\) is conserved.
The energy split becomes a probability problem.}
\label{fig:sm1-connected-systems}
\end{figure}

Given a total energy, we may ask for the probability that subsystem \(1\)
has \(E_1\) while subsystem \(2\) has the remainder. Equal boxes should have
equal energy on average, but the average is not the whole distribution. One
way to determine probabilities for a subsystem is therefore to embed it in a
larger closed system whose conservation law supplies the missing information.
This construction will recur throughout statistical mechanics.

\section{Good Laws, Bad Laws, and the Minus First Law}

Not every deterministic update rule is admissible as a fundamental law.
Consider the map
\begin{equation}
 \begin{aligned}
 R&\mapsto R,& B&\mapsto R,& Y&\mapsto R,\\
 G&\mapsto R,& O&\mapsto R,& P&\mapsto R.
 \end{aligned}
 \label{eq:sm1-all-to-red}
\end{equation}
It predicts the future perfectly: after one step the state is red. It fails
in the opposite direction. Once the system is red, nothing reveals whether
it came from blue, yellow, green, orange, purple, or red itself.

\begin{figure}[htbp]
\centering
\begin{tikzpicture}[
  >=Latex,
  state/.style={draw,circle,minimum size=8mm,inner sep=1pt}
]
\node[state,fill=red!12] (R) at (0,0) {$R$};
\node[state] (B) at (-3,1.6) {$B$};
\node[state] (Y) at (-3,0) {$Y$};
\node[state] (G) at (-3,-1.6) {$G$};
\node[state] (O) at (3,1.6) {$O$};
\node[state] (P) at (3,-1.6) {$P$};
\draw[->] (B) -- (R);
\draw[->] (Y) -- (R);
\draw[->] (G) -- (R);
\draw[->] (O) -- (R);
\draw[->] (P) -- (R);
\draw[->,loop above] (R) to (R);
\end{tikzpicture}
\caption{A many-to-one law predicts the next state but erases distinctions
among possible past states.}
\label{fig:sm1-bad-law}
\end{figure}

By contrast, a closed cycle preserves the record. If the elapsed time is
known, we can evolve the state forward or backward, even after many thousands
of periods. The law preserves distinctions between initial conditions.
Susskind calls this requirement the \emph{minus first law of physics}:
information is not fundamentally lost.

For a finite discrete state space, the diagrammatic criterion is simple:
\begin{quote}
Every state of a good law has exactly one outgoing arrow and exactly one
incoming arrow.
\end{quote}
The outgoing arrow determines the future; the incoming arrow reconstructs
the past. Mathematically, the update is a permutation of the states.

\subsection{Friction is only an apparent counterexample}

The continuous counterpart is Liouville's theorem. Friction appears to
contradict it. Slide an eraser across a table and, wherever the eraser starts,
it eventually comes to rest. A schematic drag equation is
\begin{equation}
 m\ddot x=-\gamma\dot x,
 \qquad \gamma>0.
 \label{eq:sm1-drag}
\end{equation}
Its velocity decays exponentially. If this were the rock-bottom law of every
particle in a sealed room, all motion would die away, the temperature would
fall to zero, and many initial states would collapse onto one final state.
Energy conservation would fail, information would disappear, and disorder
would decrease. This would be the continuum version of everything going to
red.

Real friction does not do this. The eraser heats the table, itself, and the
surrounding environment. Microscopic positions and momenta record the
distinctions lost from the visible motion. Friction is an effective equation
for a subsystem after unobserved degrees of freedom have been omitted; it is
not a fundamental dissipative law for the complete closed system.

\begin{classroomqa}
\classquestion{What does the differential equation
\(m\ddot x=-\gamma\dot x\) represent, and why is it not a fundamental
closed-system law?}
\lecturerresponse{It represents viscous drag. By itself it drives every
velocity toward zero, which would make a gas cool to zero regardless of its
initial temperature. In an actual closed description, the lost mechanical
energy and information move into microscopic degrees of freedom. The drag
equation is reliable only after those variables have been suppressed.
\lecturetimestamp{00:44:40--00:48:56}}
\end{classroomqa}

\section{Occupied States and the First Entropy}

Information conservation has an equivalent probabilistic statement. Suppose
\(M\) of \(N\) states are assigned equal probability,
\begin{equation}
 p_i=
 \begin{cases}
 1/M, & i\text{ in the occupied set},\\
 0, & i\text{ outside it}.
 \end{cases}
 \label{eq:sm1-uniform-occupied}
\end{equation}
An information-preserving update may reshuffle which states are occupied,
but it cannot merge them. At every later time there are still \(M\) states of
probability \(1/M\). The bad map \eqref{eq:sm1-all-to-red}, by contrast,
can collapse many occupied states into one.

The number \(M\) measures ignorance. If \(M=N\), every state is possible and
our ignorance is maximal. If \(M=1\), the microscopic state is known exactly.
The first definition of entropy is therefore
\begin{equation}
 S=\log M.
 \label{eq:sm1-entropy-logm}
\end{equation}
It gives
\begin{equation}
 S_{\max}=\log N,
 \qquad
 S_{\min}=0.
\end{equation}
If the state space is infinite, there need not be a finite upper bound.

Entropy appears before temperature in this development, and even before the
formal statement of energy conservation. In that sense it is more primitive.
Yet it is not a property of the mechanical system alone. It depends jointly
on the system and on the probability distribution representing what is
known. The same mechanical system may be assigned different entropies by
observers with different information.

Under exact microscopic evolution, the occupied states are only rearranged,
so this fine-grained entropy is conserved. In practice we lose timing,
discard correlations, and fail to track enormous numbers of variables. Our
probability distribution then spreads over more alternatives, and the
entropy assigned to the coarse description increases. The equations need
not destroy information for our usable description to lose it.

\begin{classroomqa}
\classquestion{What is the greatest possible entropy of a finite system
with \(N\) states?}
\lecturerresponse{It is \(\log N\), attained when all \(N\) states are
equally probable. The least entropy is zero, attained when one state has
probability one. Entropy therefore measures the state of knowledge as well
as the available state space.
\lecturetimestamp{00:56:46--00:58:19}}
\end{classroomqa}

\section{Liouville Flow in Continuous Phase Space}

For particles with continuous positions and momenta, a microscopic state is
a point in phase space. We write the coordinates as \(q_a\) and their
conjugate momenta as \(p_a\). With many particles there are enormously many
axes; a two-dimensional drawing with one \(q\) and one \(p\) is only a
stand-in. Phase space is even-dimensional because every coordinate has a
momentum partner.

Suppose our knowledge is represented by a patch \(\Omega_0\): the probability
density is uniform inside the patch and zero outside. For example, knowing
that every particle lies in the room restricts the position coordinates,
while bounds on their speeds restrict the momentum coordinates.

Hamiltonian evolution carries each phase-space point along a trajectory.
Nearby points move to nearby points, and the occupied patch flows like a
parcel of incompressible fluid:
\begin{figure}[htbp]
\centering
\begin{tikzpicture}[>=Latex,scale=0.95]
\draw[->] (-4.2,0) -- (4.3,0) node[right] {$q$};
\draw[->] (-3.7,-2.2) -- (-3.7,2.4) node[above] {$p$};
\draw[fill=blue!12,thick]
  (-3.0,-0.5) .. controls (-2.7,1.0) and (-1.3,1.2) .. (-0.9,0.1)
  .. controls (-0.6,-1.0) and (-2.2,-1.4) .. cycle;
\node at (-1.9,0.0) {$\Omega_0$};
\draw[->,very thick] (-0.3,0.2) -- (0.8,0.2);
\draw[fill=orange!15,thick]
  (1.2,-1.1) .. controls (1.5,0.3) and (2.3,1.8) .. (3.2,1.1)
  .. controls (3.6,0.7) and (2.8,-0.4) .. (2.5,-1.5)
  .. controls (2.1,-1.8) and (1.5,-1.6) .. cycle;
\node at (2.35,0.1) {$\Omega_t$};
\end{tikzpicture}
\caption{Hamiltonian flow may stretch and shear an occupied phase-space
region, but it preserves its volume.}
\label{fig:sm1-liouville-flow}
\end{figure}

Liouville's theorem states
\begin{equation}
 \operatorname{Vol}(\Omega_t)
 =
 \operatorname{Vol}(\Omega_0).
 \label{eq:sm1-liouville-volume}
\end{equation}
If the initial density is uniform on \(\Omega_0\), the transported density is
uniform on \(\Omega_t\). This is the continuous analog of preserving the
number and weights of occupied discrete states.

The theorem also resolves the friction puzzle geometrically. Looking only at
the eraser's visible momentum, a cloud of possibilities appears to squeeze
toward \(p=0\). In the full phase space, the cloud spreads into the positions
and momenta of molecules in the eraser and table. A projection can contract
while the full high-dimensional volume remains constant.

\begin{classroomqa}
\classquestion{If every visible trajectory approaches rest, has the
occupied phase-space region collapsed onto the axis \(p=0\)?}
\lecturerresponse{Only in the reduced projection. The complete system
contains the microscopic variables of the table, eraser, and surroundings.
As the visible momentum range narrows, the distribution spreads through
those hidden directions. Liouville's theorem applies to the full closed
phase space, not to an incomplete set of macroscopic variables.
\lecturetimestamp{01:04:37--01:06:29}}
\end{classroomqa}

\section{From the Minus First Law to the First Law}

\subsection{A brief look at the zeroth law}

Before stating the first law, it is useful to name the zeroth. Thermal
equilibrium has not yet been defined, so for now the law is only a structural
statement:
\begin{quote}
If system \(A\) is in thermal equilibrium with \(B\), and \(B\) is in thermal
equilibrium with \(C\), then \(A\) is in thermal equilibrium with \(C\).
\end{quote}
This transitivity will later make temperature possible. For the moment we set
it aside.

\subsection{Energy conservation}

The first law for a closed system is simply
\begin{equation}
 \frac{dE}{dt}=0.
 \label{eq:sm1-first-law}
\end{equation}
Its compact form conceals its power. Energy is a conserved quantity, and the
word energy will acquire more thermodynamic content later.

If a closed system consists of two parts and their energies are additive,
\begin{equation}
 E=E_1+E_2,
\end{equation}
then
\begin{equation}
 \frac{dE_1}{dt}+\frac{dE_2}{dt}=0,
 \qquad\text{or}\qquad
 \frac{dE_1}{dt}=-\frac{dE_2}{dt}.
 \label{eq:sm1-energy-exchange}
\end{equation}
The second form makes the exchange explicit: energy lost by one subsystem is
gained by the other.

The additivity assumption requires care. Two gravitating particles have
individual kinetic energies, but their interaction contributes a potential
energy belonging to the pair. The total energy is conserved even though it
cannot be written exactly as one body's energy plus the other's.

Thermodynamics often permits an excellent additive approximation because
bulk energy scales with volume while short-range interaction energy between
large neighboring pieces scales mainly with their contact area. Divide a
large table into blocks: the internal energy of each block is a volume
effect, while the coupling between adjacent blocks is a surface effect. For
large blocks, the latter is comparatively small. Equation
\eqref{eq:sm1-first-law} is exact for the closed whole; the subsystem form
\eqref{eq:sm1-energy-exchange} assumes negligible interaction energy.

\begin{classroomqa}
\audiencequestion{Does the normalization \(\sum_i p_i=1\) also assume that
the listed outcomes are mutually exclusive?}
\lecturerresponse{Yes. A state label is defined so that the alternatives do
not overlap: if the die is in the yellow state, it is not simultaneously in
the red state. Colors can be mixed in ordinary language, but the state
partition used in the probability calculation must consist of mutually
exclusive alternatives.
\lecturetimestamp{01:12:14--01:12:55}}
\end{classroomqa}

\section{Returning to Entropy}

The lecture now returns explicitly to unfinished business. Temperature,
although familiar to the senses, is a derived quantity. Entropy and energy
must be in place before temperature can be defined cleanly.

\subsection{General probability distributions}

The special definition \(S=\log M\) applies when \(M\) states are equally
probable and all others have probability zero. A general distribution need
not be rectangular. Its weights may vary smoothly or irregularly across the
state labels.

\begin{figure}[htbp]
\centering
\includegraphics[width=0.92\textwidth]{lecture_01_figure_05.png}
\caption{A probability distribution drawn over discrete state labels. The
spread, rather than the precise curve, is the feature needed for entropy.}
\lectureframe{01:14:40}
\end{figure}

\begin{figure}[htbp]
\centering
\resizebox{\linewidth}{!}{%
\begin{tikzpicture}[>=Latex,scale=0.93]
\draw[->] (-4.2,0) -- (4.4,0) node[right] {state label \(i\)};
\draw[->] (-4,0) -- (-4,3.2) node[above] {$p_i$};
\foreach \x in {-3.6,-3.1,...,3.6}
  \draw (\x,0.08) -- (\x,-0.08);
\draw[very thick]
  plot[smooth] coordinates {
  (-3.7,0.12) (-3.0,0.28) (-2.3,0.72) (-1.6,1.45)
  (-0.9,2.25) (-0.2,2.65) (0.5,2.28) (1.2,1.55)
  (1.9,0.88) (2.7,0.38) (3.5,0.14)};
\draw (-1.8,-0.18) -- (-1.8,-0.5);
\draw (1.8,-0.18) -- (1.8,-0.5);
\draw[<->] (-1.8,-0.46) -- (1.8,-0.46);
\node at (0,-0.82) {effective spread};
\end{tikzpicture}
}
\caption{A narrow distribution carries less uncertainty than a broad one.
The entropy must measure the distribution as a whole.}
\label{fig:sm1-probability-spread}
\end{figure}

The general definition is
\begin{equation}
 S=-\sum_i p_i\log p_i.
 \label{eq:sm1-shannon-entropy}
\end{equation}
Using \eqref{eq:sm1-expectation}, this is
\begin{equation}
 S=-\langle\log p\rangle.
\end{equation}
The minus sign compensates for the fact that \(\log p_i<0\) when
\(0<p_i<1\).

\subsection{Recovering the Uniform Result}

We must check that the general formula reproduces the first definition.
For an unoccupied state, \(p_i=0\) while \(\log p_i\) diverges. The product
has the well-defined limiting value
\begin{equation}
 \lim_{p\to0^+}p\log p=0.
 \label{eq:sm1-plogp-limit}
\end{equation}
Thus states of zero probability contribute nothing.

\begin{classroomqa}
\audiencequestion{If \(p_i=0\), does the term really vanish? The logarithm
of zero is minus infinity.}
\lecturerresponse{The term is defined by the limit
\(\lim_{p\to0^+}p\log p\), which is zero. The factor \(p\) approaches zero
faster than the logarithm diverges. Hence arbitrarily improbable states make
arbitrarily small contributions, and an exactly impossible state contributes
none.
\lecturetimestamp{01:18:27--01:19:46}}
\end{classroomqa}

For each of the \(M\) occupied states, \(p_i=1/M\). Therefore
\begin{align}
 S
 &=-\sum_{\mathrm{occupied}\ i}
 \frac1M\log\!\left(\frac1M\right)\\
 &=-M\frac1M\log\!\left(\frac1M\right)\\
 &=-\log\!\left(\frac1M\right)
 =\log M.
 \label{eq:sm1-entropy-check}
\end{align}
The general expression agrees with the special one exactly. A broader
distribution gives appreciable weight to more alternatives and therefore
has larger entropy; a narrower distribution has smaller entropy.

\begin{figure}[htbp]
\centering
\includegraphics[width=0.92\textwidth]{lecture_01_figure_06.png}
\caption{The marked width \(M\) emphasizes the connection between the spread
of significant probabilities and the entropy.}
\lectureframe{01:16:42}
\end{figure}

Entropy belongs to a probability distribution on the state space. It is not
a bare observable such as the energy or momentum of one microstate. An exact
state may be represented by a distribution concentrated at that state and
then has zero entropy, but the entropy changes when the state of knowledge
changes.

\section{Coins, Bits, and Explicit Counting}

Let \(n\) be the number of coins. The number of microscopic configurations is
\begin{equation}
 N=2^n.
 \label{eq:sm1-coin-state-count}
\end{equation}
If nothing is known about the coins, all \(2^n\) configurations are equally
probable, and
\begin{equation}
 S_{\max}=\log(2^n)=n\log2.
 \label{eq:sm1-coin-max-entropy}
\end{equation}
The entropy is additive in the number of independent coins.

With logarithms to base \(2\), a completely unknown two-state system has
entropy one bit. With natural logarithms it has entropy \(\log2\). The unit
is conventional; the underlying ordering and additivity are unchanged.

At the other extreme, if every coin is known, precisely one configuration
has probability one:
\begin{equation}
 S=0.
\end{equation}
More knowledge means less entropy.

Now suppose exactly one of the \(n\) coins is tails, all others are heads,
and the position of the tail is unknown. There are \(n\) equally likely
configurations, one for each possible position of the tail, so
\begin{equation}
 S=\log n.
 \label{eq:sm1-one-tail-entropy}
\end{equation}
This is generally not an integer multiple of \(\log2\). A bit is a convenient
unit, not a claim that every entropy occurs in whole-bit steps.

\begin{classroomqa}
\audiencequestion{Does \(\log2\) appear only because each coin has two
possible states? Would a three-state element give \(\log3\)?}
\lecturerresponse{Exactly. A uniformly unknown element with \(k\) possible
states contributes \(\log k\). Binary switches are especially convenient in
computers because physical switches naturally realize two stable states and
because two is the smallest nontrivial integer.
\lecturetimestamp{01:29:55--01:30:48}}
\end{classroomqa}

\begin{classroomqa}
\audiencequestion{What base is intended for the logarithm?}
\lecturerresponse{In these physics formulas, \(\log\) means the natural
logarithm. Information theorists often use base \(2\), so entropy is measured
directly in bits. Changing the base multiplies every entropy by one fixed
conversion factor and changes no physical conclusion.
\lecturetimestamp{01:30:49--01:32:06}}
\end{classroomqa}

\section{Entropy in Continuous Phase Space}

Return to a classical phase space \(\Gamma\). If the probability density is
uniform on a region of volume \(V_\Gamma\) and zero outside, the analog of
counting \(M\) equally likely states is
\begin{equation}
 S=\log\!\left(\frac{V_\Gamma}{V_0}\right).
 \label{eq:sm1-continuous-uniform}
\end{equation}
Here \(V_\Gamma\) is volume in the full high-dimensional space of coordinates
and momenta, not ordinary spatial volume. The fixed volume \(V_0\) is a
reference phase-space cell that makes the logarithm dimensionless. Setting
\(V_0=1\) gives the shorthand \(S=\log V_\Gamma\) used at the board.

For a general probability density \(P(q,p)\), the discrete sum becomes a
phase-space integral. Normalization is
\begin{equation}
 \int_\Gamma d\Gamma\,P(q,p)=1,
 \label{eq:sm1-continuous-normalization}
\end{equation}
and the entropy is
\begin{equation}
 S=-\int_\Gamma d\Gamma\,
 P(q,p)\log\!\bigl[P(q,p)V_0\bigr].
 \label{eq:sm1-continuous-entropy}
\end{equation}
The density is probability per phase-space cell. Conceptually,
\eqref{eq:sm1-continuous-entropy} still measures how broadly the probability
is distributed; for a uniform blob it reduces to the logarithm of the
occupied volume measured in units of the reference cell. Changing that fixed
cell shifts every continuous entropy by the same additive constant and leaves
entropy differences unchanged.

\begin{classroomqa}
\audiencequestion{Why did the coordinate symbol change from \(x\) to \(q\)?}
\lecturerresponse{There is no physical change. Classical mechanics commonly
uses either \(x\) or \(q\) for coordinates and \(p\) for conjugate momenta.
The switch of notation was only a slip; the phase-space construction is the
same.
\lecturetimestamp{01:37:10--01:37:39}}
\end{classroomqa}

Temperature and the Boltzmann distribution are deliberately postponed. The
next task will be to define thermal equilibrium and then identify its
probability distribution.

\section{Questions on Correlation, Additivity, and Information}

\subsection{Correlation means that one measurement changes another
distribution}

Complete ignorance about independent coins assigns equal probability to
every configuration. Measuring one coin then tells us nothing about any
other coin. Their distributions are unchanged.

The one-tail ensemble is different. If one measured coin is tails, every
other coin is certainly heads. If the measured coin is heads, the unique tail
must lie among the remaining \(n-1\) coins, so each remaining position has
conditional probability
\begin{equation}
\begin{aligned}
 p(&\text{tail at a specified remaining position}\\
   &\mid\text{measured head})=\frac{1}{n-1}.
\end{aligned}
 \label{eq:sm1-conditional-tail}
\end{equation}
This is correlation: learning one variable changes the probability
distribution assigned to another.

\begin{classroomqa}
\audiencequestion{Was independence between the coins assumed in the entropy
examples?}
\lecturerresponse{Only when complete ignorance was represented by equal
probabilities for all \(2^n\) configurations. In that case, measuring one
coin leaves the distributions of the others unchanged. The one-tail example
is explicitly correlated: a measurement updates what is possible elsewhere.
The stated information determines which probability distribution is being
used.
\lecturetimestamp{01:38:43--01:41:58}}
\end{classroomqa}

The disconnected-cycle example has the same logical form. Measuring the
conserved \(+1\) or \(-1\) label restricts all later states to one sector.
Conditioning on conserved information changes the admissible distribution.

Entropy is additive for uncorrelated systems. If
\begin{equation}
 p_{ij}=p_i^{(A)}p_j^{(B)},
\end{equation}
then
\begin{align}
 S_{AB}
 &=-\sum_{i,j}p_i^{(A)}p_j^{(B)}
 \log\!\left[p_i^{(A)}p_j^{(B)}\right]\\
 &=S_A+S_B.
 \label{eq:sm1-entropy-additivity}
\end{align}
Correlations prevent this factorization and reduce the independent
uncertainty of the pair.

As an exercise, imagine a row of binary variables for which an up-state makes
each neighbor three-quarters likely to be down. Measuring one site
immediately changes what is known about nearby sites. Such a distribution is
correlated, and its entropy cannot be found by adding isolated one-site
entropies without accounting for those correlations.

\subsection{Boltzmann, Shannon, and quantum uncertainty}

\begin{classroomqa}
\audiencequestion{How is the older thermodynamic entropy related to
Shannon's information-theoretic entropy?}
\lecturerresponse{The same functional form appears in both:
\(-\sum_i p_i\log p_i\). Shannon commonly used base \(2\), while physicists
usually use natural logarithms and restore Boltzmann's constant when entropy
is expressed in thermodynamic units. The mathematical difference is only a
fixed conversion factor, although the physical questions supplying the
probabilities can be very different.
\lecturetimestamp{01:44:05--01:46:06}}
\end{classroomqa}

For equally likely microstates, the thermodynamic formula is often written
\begin{equation}
 S_{\mathrm{thermo}}=k_{\mathrm B}\log W,
 \label{eq:sm1-boltzmann-formula}
\end{equation}
where \(W\) is the number of accessible microstates. In units with
\(k_{\mathrm B}=1\), this is precisely \(S=\log W\). Boltzmann's constant
converts between temperature units and energy units; the natural microscopic
scale of temperature is energy.

\begin{classroomqa}
\audiencequestion{Is this entropy simply another form of Heisenberg's
uncertainty principle?}
\lecturerresponse{No. The present entropy measures uncertainty in a
probability distribution over possible states. In quantum language it is the
uncertainty of a mixed state, not the intrinsic measurement uncertainty that
can remain even for a pure state. The two notions must be kept distinct.
\lecturetimestamp{01:46:07--01:47:31}}
\end{classroomqa}

\section{What Has Been Established}

The starting point was not randomness in the laws. Deterministic dynamics,
combined with incomplete knowledge of the initial state, already produces
probabilities. A single orbit assigns weights through time spent in each
state; disconnected orbits reveal conserved quantities and require prior
weights between sectors. Admissible microscopic laws preserve distinctions,
which in a finite state space means one incoming and one outgoing arrow at
every state and in continuous mechanics becomes Liouville's theorem.

That structure leads first to
\begin{equation}
 S=\log M
\end{equation}
for \(M\) equally probable occupied states, then to
\begin{equation}
 S=-\sum_i p_i\log p_i,
\end{equation}
and finally to its phase-space integral. Entropy records a probability
distribution and therefore a state of knowledge. Fine-grained information
is preserved by the microscopic law even when a practical, coarse-grained
description loses track of it.

Energy conservation supplies the first law, while additivity of subsystem
energies requires negligible interaction energy. Entropy is additive under
a parallel condition: the subsystem probabilities must be uncorrelated.
These two kinds of additivity prepare the next step, where thermal
equilibrium, temperature, and the Boltzmann distribution will connect energy
to probability.
